\documentclass{article}
\usepackage{amsmath}
\usepackage[ruled,vlined,linesnumbered]{algorithm2e}

\begin{document}

\section*{Ricerca binaria (\textit{Binary Search})}

  \vspace{0.3cm}

  \SetKwProg{Fn}{function}{}{}
  \SetKw{Return}{\textbf{return}}

  \begin{algorithm}[H]
  \caption{Search the value v among the n ordered elements of S}
  \Fn{BINARYSEARCH($S$, $v$, $\mathit{first}$, $\mathit{last}$)}{
      \If{$\mathit{first} > \mathit{last}$}{
          \Return $0$\;
      }
      $m \leftarrow \lfloor \frac{\mathit{first} + \mathit{last}}{2} \rfloor$\;
      \eIf{$S[m] = v$}{
          \Return $m$\;
      }{
          \eIf{$S[m] > v$}{
              \Return $\mathrm{binarySearch}(S, v, \mathit{first}, m - 1)$\;
          }{
              \Return $\mathrm{binarySearch}(S, v, m + 1, \mathit{last})$\;
          }
      }
  }
  \end{algorithm}
  \vspace{0.3cm}

Il cuore dell'algoritmo sta nella "divisione" dell'array, tutto in 
maniera ricorsiva. Se seguissimo "\textit{ciclo per ciclo}" notiamo
come prendiamo sempre il valore centrale, ovvero $m$, e poi 
analizziamo il sotto vettore a "destra" di $m$, o alla sua sinistra,
a seconda del valore $v$ che cerchiamo.
\\
Attenzione, i parametri $first$ e $last$ rappresentano gli 
\textbf{indici} della nostra struttura, non gli elementi. Per questo 
motivo la condizione $first > last$ ha il solo scopo di dirci se i nostri
indici si sono "scambiati", e ciò avviene solo se \textbf{non} è stato 
trovato $v$. 
\\
\\
Prendiamo come esempio un array del tipo $[1,2,3,4,9,11]$, e poniamo l'elemento
da cercare come $v = 9$; seguiamo ogni passaggio. \\
\textbf{Prima iterazione}: ignoriamo la prima condizione per quanto detto 
nel paragrafo prima e notiamo subito che la prima istruzione sarà: 
\[ 
m = \frac{first + last}{2} \rightarrow 
m = \frac{0 + 5}{2} \rightarrow m \approx 3
\]
Seguita dalla condizione $S[3] = 9$ che sarà falsa in quanto l'elemento 
alla posizione $S[3] = 4$. \\
Entriamo ora nel punto dell'algoritmo dove 
"scendiamo" nel sottoarray generato a partie da $m$, che in questo caso
sarà il sottoarray di destra, ovvero la struttura con elementi
$[9, 11]$. Ripeteremo quindi questa iterazione usando i nuovi parametri
$(S, v = 9, first = m + 1, last)$. \\
\textbf{Seconda iterazione}: ignoriamo di nuovo la prima condizione, e 
calcoliamo il nuovo $m$:
\[ 
m = \frac{first + last}{2} \rightarrow 
m = \frac{4 + 5}{2} \rightarrow m \approx 4
\]
E in questo caso $S[m] = v$ restituisce $true$ quindi restituiremo 
l'indice appena trovato $ m = 4$.

\subsection*{Costo computazionale}
Il caso migliore di questo algoritmo è ovviamente un $O(1)$ poichè se 
l'elemento da trovare ($v$) è esattamente al centro $S[m] = v$ darà 
subito $true$. \\
Il caso peggiore invece è un pochino più complesso. Innanzitutto questo caso
non è altro che lo scenario dove \textbf{non troviamo} l'elemento $v$ o quello
in cui $v$ è posto come estremo dell'array (primo o ultimo elemento). \\
Quello che farà il nostro algoritmo sarà quindi dimezzare l'array ad ogni 
iterazione:
\begin{itemize}
  \item $input_1 = n/2$
  \item $input_2 = n/2^2$
  \item $input_3 = n/2^3$
  \item $input_i = n/2^i$
\end{itemize}
Nel caso peggiore in cui l'elemento è agli estremi dell'array arriveremo 
all'ultima iterazione dove il sottoarray sarà di un solo elemento, vale a dire
che avremo un input del tipo: 
\[ n/2^i = 1 \]
Adesso ci basta solo applicare della matematica per trovare che:
\[ n = 2^i \rightarrow i = \log_2 n\]
Il caso peggiore del \textbf{\textit{binary search}}
è asintoticamente pari a $O(\log n)$


\end{document}